% ATENÇÃO (nota de trabalho, não sai no PDF): capítulo vivo.
% Substituiu integralmente o conteúdo de exemplo do template (lipsum + \Gls de
% um TCC de outra área). Estrutura por eixo de proximidade, não por autor.

\chapter{Trabalhos Relacionados}
\label{cap:trabalhos-relacionados}

Este capítulo situa o trabalho em relação a quatro linhas de pesquisa que o tangenciam: geração de modelos de processo por modelos de linguagem, representações compactas de processos, compressão de \textit{tokens} em sistemas com LLMs, e garantia de validade sintática na geração de linguagens formais. O critério de organização é a proximidade com a pergunta desta pesquisa — o custo de gerar marcação hierárquica verbosa —, e não a cronologia ou a autoria.

\section{Geração de Modelos de Processo com LLMs}
\label{sec:tr-geracao}

A referência mais próxima em objetivo é o ProMoAI \cite{kourani2024promoai}, que emprega um LLM de fronteira, guiado por \textit{prompt}, para produzir código POWL executável, posteriormente traduzido para BPMN. A estratégia compartilha com este trabalho a premissa central de \textbf{não} pedir XML diretamente ao modelo, interpondo uma representação intermediária mais adequada às capacidades do gerador.

As diferenças são de natureza e de escopo. O ProMoAI depende de um modelo de fronteira em tempo de inferência, o que mantém o custo por geração no patamar que este trabalho procura reduzir; a representação intermediária é código Python, escolhida por familiaridade do modelo com tarefas de programação, e não projetada para minimizar \textit{tokens} emitidos; e a compressão obtida não é medida nem tratada como objetivo. Aqui, ao contrário, a compressão é métrica de primeira classe, a representação é projetada para isso, e o gerador-alvo é um modelo pequeno especializado por ajuste supervisionado. A reexecução do ProMoAI está fora do escopo desta versão, de modo que ele é tratado como referência qualitativa; comparar com números publicados sob protocolo distinto seria metodologicamente indefensável.

\section{Representações Compactas de Processos}
\label{sec:tr-representacoes}

Diversas notações alternativas ao XML BPMN foram propostas com objetivos que se sobrepõem parcialmente ao deste trabalho. A representação \textit{JSON branches} \cite{kopke2024efficient} foi introduzida explicitamente para reduzir contagem de \textit{tokens} e favorecer aderência a esquema na geração por LLMs. A representação \textit{Process Model Elements} \cite{volter2024generative} organiza o modelo em listas planas de elementos, facilitando contagem e comparação. Somam-se a elas notações de visualização como Mermaid e Graphviz, e simplificações do próprio XML.

O trabalho de \citeonline{brissard2025pmr} é o que mais se aproxima em objeto empírico: realiza análise comparativa de nove representações de modelos de processo para geração com LLMs, sobre o mesmo conjunto de dados aqui utilizado como \textit{holdout}. É, por isso, o antecedente que exige posicionamento mais cuidadoso.

A distinção é de \textbf{pergunta de pesquisa}, e não de artefato. Aquele trabalho pergunta qual representação produz os melhores modelos de processo quando um LLM é instruído por \textit{prompt} — uma pergunta sobre adequação de notação à tarefa de modelagem. Esta pesquisa pergunta se a estratégia de comprimir a saída e expandi-la deterministicamente torna economicamente viável a geração de marcação hierárquica verbosa por modelos pequenos, usando BPMN como caso de estudo instrumentado. Disso decorrem três diferenças operacionais:

\begin{alineas}
	\item \textbf{o gerador}: representações avaliadas por \textit{prompt} em modelos de fronteira, contra um modelo pequeno especializado por ajuste supervisionado, cujo custo por inferência é ordens de grandeza menor;
	\item \textbf{a garantia sintática}: notações comparadas quanto à qualidade do modelo produzido, contra uma cadeia em que a validade perante o esquema é assegurada por construção, deixando de ser variável aleatória;
	\item \textbf{a métrica de custo}: compressão medida e reportada como resultado de primeira ordem, e não como propriedade qualitativa da notação.
\end{alineas}

Que representações compactas de BPMN já existam não é objeção a este trabalho — é evidência de que o problema é reconhecido. A contribuição pretendida não é a notação em si, mas a demonstração empírica de que a combinação entre compressão medida, expansão determinística verificável e especialização de modelo pequeno desloca a fronteira de custo da tarefa. Registre-se ainda que a DSL proposta cobre construções que algumas dessas notações não representam — múltiplos eventos de início e fim, fluxos de mensagem e raias —, embora essa maior expressividade seja consequência do desenho, e não sua justificativa.

\section{Compressão de \textit{Tokens}: Entrada \textit{versus} Saída}
\label{sec:tr-compressao}

A redução do consumo de \textit{tokens} em sistemas com LLMs é problema ativo, e a literatura predominante trata do lado da \textbf{entrada}. O LLMLingua \cite{jiang2023llmlingua} comprime \textit{prompts} removendo conteúdo de baixa informação segundo um modelo de linguagem auxiliar, preservando o desempenho na tarefa a jusante. Trabalhos de gestão de contexto perseguem objetivo análogo por seleção e sumarização do que entra na janela.

Este trabalho ataca o lado oposto e complementar: o custo da \textbf{saída}. A assimetria é relevante e raramente explicitada. \textit{Tokens} de entrada podem ser processados em paralelo e beneficiam-se de mecanismos de cache; \textit{tokens} de saída são gerados sequencialmente, dominam a latência percebida e são tipicamente tarifados a preço superior. Em tarefas cuja saída é marcação verbosa, o gargalo econômico está na geração, não na leitura — e comprimir o \textit{prompt} não o alivia. As duas famílias de técnica são, portanto, ortogonais e combináveis, e não concorrentes.

\section{Modelos Pequenos Especializados}
\label{sec:tr-slm}

A viabilidade da proposta depende de que um modelo de pequeno porte, especializado, alcance desempenho competitivo em tarefa estreita. Essa premissa é sustentada por evidência recente sobre o papel de modelos pequenos em sistemas agênticos \cite{belcak2025small}, que argumenta serem eles suficientes — e economicamente preferíveis — para subtarefas repetitivas e bem delimitadas, reservando modelos de fronteira para as etapas que exigem generalidade.

A geração de marcação estruturada a partir de descrição textual é exatamente uma dessas subtarefas: formato fixo, vocabulário restrito, critério de sucesso verificável. Métodos de ajuste eficiente em parâmetros \cite{hu2022lora}, especialmente em sua variante quantizada \cite{dettmers2023qlora}, tornam a especialização exequível com recursos computacionais modestos, o que é condição para que o argumento econômico se sustente.

\section{Garantia de Validade Sintática}
\label{sec:tr-validade}

Há uma alternativa à estratégia adotada que precisa ser reconhecida: a \textbf{decodificação restrita}, que impõe a gramática do formato-alvo durante a amostragem, impedindo por construção a emissão de sequências inválidas \cite{poesia2022synchromesh}. Aplicada ao XML BPMN, garantiria validade sintática sem representação intermediária alguma.

Duas considerações delimitam a comparação. A primeira é que a decodificação restrita resolve a validade, mas \textbf{não} o custo: o modelo continua emitindo cada \textit{token} da marcação verbosa, de modo que a economia perseguida aqui permanece inalcançada — e, com efeito, as duas técnicas são combináveis, podendo a decodificação restrita ser aplicada sobre a própria DSL. A segunda é que restringir a decodificação a uma gramática garante boa formação sintática, mas não coerência estrutural: um documento pode respeitar o esquema e ainda referenciar identificadores inexistentes ou descrever um grafo desconexo, classe de erro que a expansão determinística elimina por não permitir que ela seja expressa.

A avaliação empírica da decodificação restrita como braço experimental está fora do escopo desta versão, e é registrada como trabalho futuro derivado.
